\begin{table}[H]
\centering
\caption{Approximate Peer/Rank Diagnostic (Borusyak--Hull Control Function)}
\label{tab:peer_effects}
\begin{threeparttable}
\begin{tabular}[t]{lcc}
\toprule
\multicolumn{1}{c}{ } & \multicolumn{2}{c}{Experiments} \\
\cmidrule(l{3pt}r{3pt}){2-3}
\multicolumn{1}{c}{ } & \multicolumn{1}{c}{(1)} & \multicolumn{1}{c}{(2)} \\
\cmidrule(l{3pt}r{3pt}){2-2} \cmidrule(l{3pt}r{3pt}){3-3}
 & \multicolumn{1}{c}{Kenya} & \multicolumn{1}{c}{Liberia} \\
\midrule
Peer mean EB ability &  -0.133** &  -0.048 \\
 & (  0.052) & (  0.108) \\
\addlinespace
N & 4954 & 3153 \\
\bottomrule
\end{tabular}
\begin{tablenotes}[para]
\item Notes: Dependent variable is standardized endline score. Peer mean EB ability is the mean of classmates' EB ability, excluding the focal student. This is the approximate control function implementation: controls include baseline-decile $\times$ treatment $\times$ grade FE and strata FE\@. Standard errors are clustered at the school level in Kenya and at the school grade group level in Liberia. ***, **, and * indicate significance at 1\%, 5\%, and 10\%.
\end{tablenotes}
\end{threeparttable}
\end{table}
